%%% FIELD gabriel-scope-statement
Results 1 and 3 of Gabriel, Sachs, and Sj\"olander study their specified low-cardinality coarsening configurations with one off-target level.  Their structural formulations take \(X=g_X(X^*)\), retain a common target response \(Y_t=f_Y(U,t,\epsilon_y)\) for each well-defined \(t\in\mathcal T\), and represent the off-target outcome either by the realized fine version \(g_Y(U,X_z^*,\epsilon_y)\) in the ill-definition formulation or by the arm-specific function \(f_{Y^2}(U,z,\epsilon_y)\) in the contamination formulation.  The two results report bounds for specified target contrasts under the relevance conditions stated there.
%%% FIELD gabriel-transfer-statement
In each one-off-target NPSEM described by the cited Gabriel scope lemma, the map
\[
 X_z=g_X(X_z^*),\qquad
 Y_t=f_Y(U,t,\epsilon_y),
\]
together with
\[
 R_{c,z}=g_Y(U,X_z^*,\epsilon_y)
 \quad\text{or}\quad
 R_{c,z}=f_{Y^2}(U,z,\epsilon_y)
\]
for, respectively, the ill-definition or contamination formulation, sends every randomized source law satisfying \(P(Z=z)=\lambda_z>0\) for all \(z\in\mathcal Z\) to a law in \(\mathcal M_{\mathcal T,\{c\}}\) with exactly the same observed arm law \(p\) and target law \(\operatorname{law}(Y_{\mathcal T})\).  Conversely, every full response-type law on \(\mathcal R_{\mathrm F}\) with \(\mathcal C=\{c\}\) has a representation of each of these two forms using the full response type as \(U\) and degenerate independent errors.  The representations preserve randomization, \(p\), and \(\operatorname{law}(Y_{\mathcal T})\).  Consequently the source and present response-law fibers, and hence their target-law images, agree at the matching cardinalities.  This equality remains true after restricting both sides to any relevance condition that is a restriction on the observed law \(p\).
%%% FIELD gabriel-transfer-proof
We prove both directions constructively.  First consider a source NPSEM satisfying \(P(Z=z)=\lambda_z>0\) for every \(z\in\mathcal Z\), and collect all of its pre-assignment latent variables and errors into a random element \(L\).  Random assignment means \(Z\perp\!\!\!\perp L\).  For every \(z\in\mathcal Z\), let \(X_z^*\) be the fine potential exposure generated by the source exposure equation and put
\[
X_z=g_X(X_z^*).
\]
For every target level \(t\in\mathcal T\), define
\[
Y_t=f_Y(U,t,\epsilon_y).
\]
The right-hand side has no \(z\)-argument, so the target-exclusion condition holds.  In the ill-definition formulation set
\[
R_{c,z}=g_Y(U,X_z^*,\epsilon_y),
\]
and in the contamination formulation set
\[
R_{c,z}=f_{Y^2}(U,z,\epsilon_y).
\]
All coordinates of \((X_{\mathcal Z},Y_{\mathcal T},R_{c,\mathcal Z})\) are measurable functions of \(L\).  Since independence is preserved by measurable maps, randomization gives
\[
Z\perp\!\!\!\perp(X_{\mathcal Z},Y_{\mathcal T},R_{c,\mathcal Z}).
\]
Set \(X=X_Z\).  On \(\{X=t\}\) use the source target equation, and on \(\{X=c,Z=z\}\) use the displayed off-target equation.  This verifies exposure consistency, target-outcome consistency, and off-target-outcome consistency one at a time.  Hence the induced law belongs to \(\mathcal M_{\mathcal T,\{c\}}\).

Because \(P(Z=z)=\lambda_z>0\), conditioning on arm \(z\) is well defined.  Randomization and the consistency equations imply, for \(y\in\mathcal Y\),
\[
\begin{aligned}
P(X=t,Y=y\mid Z=z)
 &=P(X_z=t,Y_t=y), &&t\in\mathcal T,\\
P(X=c,Y=y\mid Z=z)
 &=P(X_z=c,R_{c,z}=y).
\end{aligned}
\]
These are the same events and the same probabilities as in the source NPSEM by construction.  Thus every coordinate of \(p\) is preserved.  The random vector \(Y_{\mathcal T}\) was retained coordinatewise, so its entire joint law is also preserved.

For the converse, let \(F=(x_{\mathcal Z},y_{\mathcal T},r_{c,\mathcal Z})\) have an arbitrary probability law \(\pi\) on the full response-type space with \(\mathcal C=\{c\}\).  Take \(U=F\), let every auxiliary exposure and outcome error be a constant, and draw \(Z\) independently of \(U\) with the given positive probabilities \((\lambda_z)_z\).  Constant errors are mutually independent and independent of \((U,Z)\).  Define
\[
f_X(U,z)=x_z(U),
\qquad
f_Y(U,t,0)=y_t(U).
\]
For the contamination representation, take the fine and recorded exposure labels to coincide, let \(g_X\) be the identity, and define
\[
f_{Y^2}(U,z,0)=r_{c,z}(U).
\]
Then the resulting vector of response coordinates is exactly \(F\), not merely equal to it in distribution.

For the ill-definition representation, introduce the finite fine support
\[
\mathcal X^*=\mathcal T\mathbin{\dot\cup}
\{(c,z):z\in\mathcal Z\},
\]
and the deterministic coarsening
\[
g_X(t)=t\quad(t\in\mathcal T),
\qquad g_X(c,z)=c.
\]
For each arm define
\[
X_z^*(U)=
\begin{cases}
x_z(U),&x_z(U)\in\mathcal T,\\
(c,z),&x_z(U)=c,
\end{cases}
\]
and define the outcome function on all fine versions by
\[
g_Y(U,t,0)=y_t(U),
\qquad
g_Y(U,(c,z),0)=r_{c,z}(U).
\]
It follows coordinatewise that \(g_X(X_z^*)=x_z(U)\), that the common response at target \(t\) is \(y_t(U)\), and that the response on the off-target branch in arm \(z\) is \(r_{c,z}(U)\).  Thus this representation also reproduces exactly the original full response type.  Drawing \(Z\) independently and applying consistency therefore reproduces the same observed law \(p\), the same target law, and the same randomization relation.

The two constructions are inverse at the level relevant to response-law fibers: the forward construction records \((X_{\mathcal Z},Y_{\mathcal T},R_{c,\mathcal Z})\), and the reverse construction realizes any law of precisely that vector.  Finally, let \(\mathcal P_{\mathrm{rel}}\) be any subset of the observed-law simplex.  Both maps leave \(p\) unchanged, so a law has \(p\in\mathcal P_{\mathrm{rel}}\) before the map if and only if it does afterward.  Restricting to \(\mathcal P_{\mathrm{rel}}\), including any published relevance condition that is stated solely in terms of \(p\), therefore preserves the fiber equality and its target-law image.
%%% FIELD song-scope-statement
The fully excluded finite categorical instrumental-variable model used by Song, Guo, Chan, and Richardson represents a unit by an exposure response vector \((X_z:z\in\mathcal Z)\) and one arm-invariant potential outcome \(Y_x\in\mathcal Y\) for every \(x\in\mathcal X\), together with random assignment and consistency.
%%% FIELD song-common-equivalence-statement
Let
\[
\mathcal R_{\mathrm{FE}}=\mathcal X^{\mathcal Z}\times\mathcal Y^{\mathcal X}
\]
be the fully excluded response-type space, with atoms \(g=(x_{\mathcal Z},(y_x)_{x\in\mathcal X})\).  The map
\[
J(g)=\bigl(x_{\mathcal Z},(y_t)_{t\in\mathcal T},
(r_{c,z}=y_c)_{c\in\mathcal C,z\in\mathcal Z}\bigr)
\]
is a bijection from \(\mathcal R_{\mathrm{FE}}\) onto the subset of \(\mathcal R_{\mathrm F}\) satisfying \(r_{c,z}=r_{c,z'}\) for every \(c\in\mathcal C\) and \(z,z'\in\mathcal Z\).  Its affine extension to probability laws preserves every observed cell probability and the target-law map.  Thus, after off-target potential outcomes are omitted from the reported target parameter, the fully excluded response-law class is exactly \(\mathcal M_{\mathrm{common}}\).  When \(\mathcal C\ne\varnothing\), this correspondence by itself does not identify that class with \(\mathcal M_{\mathrm{sat}}\).
%%% FIELD song-common-equivalence-proof
Define the inverse on a common-response full type
\[
f=(x_{\mathcal Z},y_{\mathcal T},r_{\mathcal C,\mathcal Z})
\quad\text{with}\quad
r_{c,z}=r_{c,z'}
\]
by retaining \(x_{\mathcal Z}\), setting \(y_t^{\mathrm{FE}}=y_t\) for \(t\in\mathcal T\), and setting
\[
y_c^{\mathrm{FE}}=r_{c,z_0}
\]
for any fixed \(z_0\in\mathcal Z\).  The common-response equalities make this definition independent of the choice of \(z_0\).  It is immediate coordinate by coordinate that this inverse composed with \(J\), and \(J\) composed with this inverse, are identity maps.  Hence \(J\) is a bijection.

For a law \(\omega\) on \(\mathcal R_{\mathrm{FE}}\), define its fully excluded observed incidence by
\[
(A_{\mathrm{FE}}\omega)_{z,(x,y)}
=\sum_{g\in\mathcal R_{\mathrm{FE}}}
\omega_g\mathbf 1\{x_z(g)=x,\ y_x(g)=y\}.
\]
Let \(\pi=J_{\#}\omega\).  If \(x=t\in\mathcal T\), then the definition of \(J\) gives \(y_t(Jg)=y_t(g)\); if \(x=c\in\mathcal C\), it gives \(r_{c,z}(Jg)=y_c(g)\).  Therefore, for every \(z,x,y\),
\[
(A_{\mathrm F}\pi)_{z,(x,y)}
=\sum_g\omega_g\mathbf 1\{x_z(g)=x,\ y_x(g)=y\}
=(A_{\mathrm{FE}}\omega)_{z,(x,y)}.
\]
Likewise, for every \(y_{\mathcal T}\in\mathcal Y^{\mathcal T}\),
\[
(B_{\mathrm F}\pi)_{y_{\mathcal T}}
=\sum_g\omega_g
\mathbf 1\{(y_t(g))_{t\in\mathcal T}=y_{\mathcal T}\},
\]
which is exactly the fully excluded target marginal after the coordinates \((Y_c:c\in\mathcal C)\) are omitted from the reported parameter.

At the causal-law level, draw either response type independently of \(Z\), set \(X=X_Z\), and use the appropriate response coordinate for \(Y\).  Randomization and the two consistency conditions then identify the displayed incidence probabilities with the observed conditional arm probabilities.  The support condition defining the image of \(J\) is precisely \(R_{c,z}=R_{c,z'}=:R_c\) for every \(c,z,z'\), which is the defining restriction of \(\mathcal M_{\mathrm{common}}\).  Surjectivity of \(J\) proves equality with that class at response-law level.  It furnishes no equality with the saturated class when \(\mathcal C\ne\varnothing\); the strictness claim requires the separate witness proved below.
%%% FIELD quotient-statement
For every \(P\in\mathcal M_{\mathcal T,\mathcal C}\) with \(\lambda_z=P(Z=z)>0\) for every \(z\in\mathcal Z\), observed conditional arm law \(p\), and skeleton \(s\), \(\mathcal I_{\mathrm F}(p)=\Theta(s)\). More precisely, \(A\kappa_{\#}\pi=s\) and \(B\kappa_{\#}\pi=B_{\mathrm F}\pi\) for every feasible full law \(\pi\); conversely, for every \(q\in F(s)\), the explicit refinement satisfies \(A_{\mathrm F}\mathfrak R_p(q)=p\), \(\kappa_{\#}\mathfrak R_p(q)=q\), and \(B_{\mathrm F}\mathfrak R_p(q)=Bq\), including every zero-\(\star\)-mass boundary. Hence the sharp set and the dimensions \(|\mathcal R|=(|\mathcal T|+1)^m2^{|\mathcal T|}\) and \(md\) are independent of \(|\mathcal C|\). This strictly extends Results 1 and 3 of Gabriel, Sachs, and Sj\"olander from their specified low-cardinality target contrasts to the complete law of \(Y_{\mathcal T}\) for arbitrary finite \(m,\mathcal T,\mathcal C\).
%%% FIELD quotient-proof
Write a full response type as
\[
f=(x_{\mathcal Z},y_{\mathcal T},r_{\mathcal C,\mathcal Z})
\in\mathcal R_{\mathrm F}.
\]
Because all potential variables have finite support, the given causal law \(P\) induces the probability vector
\[
\pi^P=\operatorname{law}_P(X_{\mathcal Z},Y_{\mathcal T},R_{\mathcal C,\mathcal Z})
\in\Delta^{|\mathcal R_{\mathrm F}|-1}.
\]
For a target cell, randomization, exposure consistency, target exclusion, and target-outcome consistency give
\[
(A_{\mathrm F}\pi^P)_{z,(t,y)}
=P(X_z=t,Y_t=y)
=P(X=t,Y=y\mid Z=z)
=p_z(t,y).
\]
For an off-target cell, randomization, exposure consistency, and off-target-outcome consistency analogously give
\[
(A_{\mathrm F}\pi^P)_{z,(c,y)}
=P(X_z=c,R_{c,z}=y)
=P(X=c,Y=y\mid Z=z)
=p_z(c,y).
\]
Every conditioning operation is valid because \(\lambda_z>0\).  Thus \(A_{\mathrm F}\pi^P=p\), so the full feasible fiber is nonempty.

We first prove the outer inclusion by a pushforward argument.  Let \(\pi\in\Delta^{|\mathcal R_{\mathrm F}|-1}\) satisfy \(A_{\mathrm F}\pi=p\), and put \(q=\kappa_{\#}\pi\).  Since a pushforward of a probability measure is a probability measure, \(q\in\Delta^{|\mathcal R|-1}\).  For \(z\in\mathcal Z\), \(t\in\mathcal T\), and \(y\in\mathcal Y\), the definitions of \(A\) and \(\kappa\) give
\[
\begin{aligned}
(Aq)_{z,(t,y)}
&=\sum_{r\in\mathcal R}q_r
  \mathbf 1\{x_z(r)=t,\ y_t(r)=y\}\\
&=\sum_{f\in\mathcal R_{\mathrm F}}\pi_f
  \mathbf 1\{x_z(f)=t,\ y_t(f)=y\}\\
&=(A_{\mathrm F}\pi)_{z,(t,y)}\\
&=p_z(t,y)=s_z(t,y).
\end{aligned}
\]
The incidence equality is algebraic.  Its causal interpretation follows from randomization, exposure consistency, target exclusion, and target-outcome consistency: positivity of \(\lambda_z\) permits conditioning and
\[
p_z(t,y)=P(X_z=t,Y_t=y).
\]
For the collapsed row, disjointness of \(\mathcal T\) and \(\mathcal C\), binary outcome support, and the definition of \(A_{\mathrm F}\) yield
\[
\begin{aligned}
(Aq)_{z,\star}
&=\sum_f\pi_f\mathbf 1\{x_z(f)\in\mathcal C\}\\
&=\sum_{c\in\mathcal C}\sum_{y\in\mathcal Y}
  (A_{\mathrm F}\pi)_{z,(c,y)}\\
&=\sum_{c\in\mathcal C}\sum_{y\in\mathcal Y}p_z(c,y)
=s_z(\star).
\end{aligned}
\]
Thus \(Aq=s\), so \(q\in F(s)\).  Moreover, \(\kappa\) retains the complete vector \(y_{\mathcal T}\).  Consequently, for every \(y_{\mathcal T}\in\mathcal Y^{\mathcal T}\),
\[
\begin{aligned}
(Bq)_{y_{\mathcal T}}
&=\sum_f\pi_f\mathbf 1\{y_{\mathcal T}(f)=y_{\mathcal T}\}\\
&=(B_{\mathrm F}\pi)_{y_{\mathcal T}}.
\end{aligned}
\]
It follows that \(\mathcal I_{\mathrm F}(p)\subseteq\Theta(s)\).

We next prove attainability by the explicit right inverse.  Fix \(q\in F(s)\).  Conditional on a reduced atom
\[
r=((x_z(r))_{z\in\mathcal Z},y_{\mathcal T}(r)),
\]
perform the refinement in \(\mathfrak R_p\): whenever \(x_z(r)=\star\), sample \((C_z,V_z)\) according to \(\eta_z\), set \(X_z=C_z\) and \(R_{C_z,z}=V_z\), and fill every unused off-target response coordinate with zero; when \(x_z(r)=t\in\mathcal T\), set \(X_z=t\) and again fill unused off-target coordinates with zero.  The draws may be taken independently across the finitely many arms conditional on \(r\).  Let
\[
\pi=\mathfrak R_p(q).
\]
Every full atom generated from \(r\) collapses back to \(r\).  Hence, by conditioning on \(r\),
\[
\kappa_{\#}\mathfrak R_p(q)=q,
\qquad
B_{\mathrm F}\mathfrak R_p(q)=Bq.
\]
For a target cell, the refinement adds no randomness relevant to the incidence indicator, and therefore
\[
\begin{aligned}
(A_{\mathrm F}\pi)_{z,(t,y)}
&=\sum_rq_r\mathbf 1\{x_z(r)=t,\ y_t(r)=y\}\\
&=(Aq)_{z,(t,y)}=s_z(t,y)=p_z(t,y).
\end{aligned}
\]
For an off-target cell, conditional independence is more than is needed; the marginal refinement at arm \(z\) gives
\[
\begin{aligned}
(A_{\mathrm F}\pi)_{z,(c,y)}
&=\sum_rq_r\mathbf 1\{x_z(r)=\star\}\eta_z(c,y)\\
&=(Aq)_{z,\star}\eta_z(c,y)\\
&=s_z(\star)\eta_z(c,y).
\end{aligned}
\]
If \(s_z(\star)>0\), the definition of \(\eta_z\) makes the last quantity \(p_z(c,y)\).  If \(s_z(\star)=0\), then nonnegativity of \(q\) and
\[
0=(Aq)_{z,\star}
=\sum_rq_r\mathbf 1\{x_z(r)=\star\}
\]
imply that every reduced atom having \(x_z(r)=\star\) has zero \(q\)-mass.  Nonnegativity of the observed cells and
\[
0=s_z(\star)=\sum_{c\in\mathcal C,y\in\mathcal Y}p_z(c,y)
\]
also imply \(p_z(c,y)=0\) for each \((c,y)\).  Thus the displayed incidence equality holds for every fixed boundary choice of \(\eta_z\).  If \(\mathcal C=\varnothing\), the sum defining \(s_z(\star)\) is zero and no positive-mass \(\star\)-branch occurs, so the singleton convention is sufficient.

It remains to verify that the constructed table is attainable in the maintained causal model.  Draw the full response type with law \(\pi\), draw \(Z\) independently with probabilities \((\lambda_z)_z\), set \(Y_{t,z}=Y_t\), put \(X=X_Z\), and define
\[
Y=\begin{cases}
Y_t,&X=t\in\mathcal T,\\
R_{c,Z},&X=c\in\mathcal C.
\end{cases}
\]
This finite construction satisfies randomization, target exclusion, exposure consistency, target-outcome consistency, and off-target-outcome consistency.  The preceding cell equations give \(A_{\mathrm F}\pi=p\).  Thus every \(Bq\in\Theta(s)\) is attained by a law in \(\mathcal M_{\mathcal T,\mathcal C}\), proving
\[
\mathcal I_{\mathrm F}(p)=\Theta(s).
\]
This is equality of the causal identified set and the observed-law functional, not a definition of the former; because every point on the right has just been attained, the equality is sharp rather than merely an outer bound.

A reduced atom has \((|\mathcal T|+1)^m\) possible exposure-response vectors and \(2^{|\mathcal T|}\) possible target-outcome vectors.  Therefore
\[
|\mathcal R|=(|\mathcal T|+1)^m2^{|\mathcal T|}.
\]
The map \(A\) records \(d=2|\mathcal T|+1\) cells in each of \(m\) arms, so its declared codomain has dimension \(md\).  Neither number involves \(|\mathcal C|\).

Finally, the claimed comparison with Gabriel, Sachs, and Sj\"olander is a model-transfer claim, not an inference from cardinalities alone.  The Gabriel NPSEM transfer lemma proves that, in each of their one-off-target configurations, the published NPSEM response-law fiber and the present saturated response-law fiber are equivalent, preserving \(p\), randomization, and the full target law.  If \(\psi\) is any of their reported target contrasts, it is a functional of \(\mu=\operatorname{law}(Y_{\mathcal T})\); hence the equality already proved gives the sharp contrast set as the computable image
\[
\{\psi(\mu):\mu\in\mathcal I_{\mathrm F}(p)\}
=\{\psi(Bq):q\in F(s)\}.
\]
The transfer remains valid after imposing any of their observed-law relevance restrictions.  The cited Gabriel scope lemma records that Results 1 and 3 cover specified low-cardinality contrasts.  The present equality instead covers every coordinate of the complete joint law for arbitrary finite \(m,\mathcal T,\mathcal C\), including configurations outside those cardinalities.  The extension is therefore strict in scope; no claim of strict numerical improvement over a published bound in a matching configuration is being made.
%%% FIELD witness-statement
For \(m\ge2\) and nonempty \(\mathcal C\), the common-response restriction \(R_{c,z}=R_c\) defines a strict submodel \(\mathcal M_{\mathrm{common}}\subsetneq\mathcal M_{\mathrm{sat}}\). For \(p^{\mathrm w}\) on arms \(\{1,2\}\), the response-type certificate \(\mathfrak C^{\mathrm w}\) exhaustively enumerates \(|\mathcal R|=36\), \(|\mathcal R_{\mathrm F}|=144\) under \(\mathcal M_{\mathrm{sat}}\), and \(72\) types under \(\mathcal M_{\mathrm{common}}\). Its feasible primal mixtures and valid endpoint inequalities certify \(\tau\in[-1/2,1]\) under \(\mathcal M_{\mathrm{sat}}\) and \(\tau\in[0,1]\) under \(\mathcal M_{\mathrm{common}}\).
%%% FIELD witness-proof
The Song common-response equivalence lemma first fixes the comparator class: the affine map \(Y_c\mapsto(R_{c,z}=Y_c)_z\), with inverse \(R_c\mapsto Y_c=R_c\), identifies the fully excluded response-law class with \(\mathcal M_{\mathrm{common}}\) after the off-target coordinates are omitted from the reported target parameter.  This is only a class equivalence; it does not equate that class's projected identified set with \(\Theta(s)\).

By definition, every common-response law is saturated, so
\[
\mathcal M_{\mathrm{common}}\subseteq\mathcal M_{\mathrm{sat}}.
\]
The inclusion is strict when \(m\ge2\) and \(\mathcal C\ne\varnothing\).  Choose \(c\in\mathcal C\) and distinct arms \(z,z'\), and take a degenerate response type with
\[
X_u=c\quad\text{for every }u\in\mathcal Z,
\qquad R_{c,z}=0,
\qquad R_{c,z'}=1.
\]
Assign arbitrary fixed values to every other response coordinate.  Draw \(Z\) independently with positive arm probabilities, set \(Y_{t,u}=Y_t\), put \(X=X_Z\), and set \(Y=R_{c,Z}\).  Randomization, target exclusion, exposure consistency, and off-target-outcome consistency follow directly; target-outcome consistency is vacuous because \(X=c\) almost surely.  The resulting full causal law belongs to \(\mathcal M_{\mathrm{sat}}\) but violates the almost-sure equality required by \(\mathcal M_{\mathrm{common}}\), proving strict inclusion.

Now specialize to the witness.  Ordering coordinates as in the certificate, direct Cartesian-product enumeration gives
\[
|\{a,b,\star\}^{\{1,2\}}\times\{0,1\}^{\{a,b\}}|
=3^2 2^2=36,
\]
\[
|\{a,b,c\}^{\{1,2\}}\times\{0,1\}^{\{a,b\}}
\times\{0,1\}^{\{c\}\times\{1,2\}}|
=3^2 2^2 2^2=144,
\]
and
\[
|\{a,b,c\}^{\{1,2\}}\times\{0,1\}^{\{a,b\}}
\times\{0,1\}^{\{c\}}|
=3^2 2^2 2=72.
\]
Each factor ranges over its entire finite support, so these lists are exhaustive.  In particular, the certificate's linear programs expose every deterministic column and their probability-simplex mixtures expose every response-type law.

The quotient theorem also supplies a stronger cardinality check.  For any finite off-target alphabet, any full observed law with skeleton \(s\), and any direction \(h\in\mathbb R^{2^{|\mathcal T|}}\), compactness of the finite simplices and \(\mathcal I_{\mathrm F}(p)=\Theta(s)\) imply
\[
\max_{\substack{\pi\ge0,\ \mathbf 1^\top\pi=1\\A_{\mathrm F}\pi=p}}
h^\top B_{\mathrm F}\pi
=
\max_{\substack{q\ge0,\ \mathbf 1^\top q=1\\Aq=s}}
h^\top Bq,
\]
and similarly for minima.  Thus checking arbitrary support directions across different finite refinements of the same skeleton must give exact equality; in the witness the reduced side uses the enumerated \(36\) columns.

For the saturated lower endpoint, place mass \(1/2\) on each of the full types
\[
(a,a,0,0,0,0),
\qquad
(c,c,1,0,1,0),
\]
where the coordinates are \((x_1,x_2,y_a,y_b,r_{c,1},r_{c,2})\).  The first atom produces \((a,0)\) in both arms; the second produces \((c,1)\) in arm \(1\) and \((c,0)\) in arm \(2\).  Their mixture therefore has observed law \(p^{\mathrm w}\), and its contrast is
\[
\tau=\tfrac12(0-0)+\tfrac12(0-1)=-\tfrac12.
\]
For the upper endpoint, place mass \(1/2\) on each of
\[
(a,a,0,1,0,0),
\qquad
(c,c,0,1,1,0).
\]
This has the same observed law and \(\tau=1\).

We certify that no saturated mixture exceeds these endpoints.  Binary support gives
\[
0\le E(Y_a),E(Y_b)\le1,
\qquad
\tau=E(Y_b)-E(Y_a)\le1.
\]
The two observed target cells and randomization plus target consistency give
\[
P(X_1=a,Y_a=0)=P(X_2=a,Y_a=0)=\tfrac12.
\]
Thus
\[
P(Y_a=0)
\ge P(X_1=a\ \text{or}\ X_2=a)
\ge\tfrac12,
\]
so \(E(Y_a)\le1/2\).  Since \(E(Y_b)\ge0\),
\[
\tau\ge-E(Y_a)\ge-\tfrac12.
\]
The two feasible mixtures attain the two valid inequalities.  The feasible response-law set is convex and \(\tau\) is linear, so mixing the two endpoint laws attains every intermediate value.  The saturated sharp interval is therefore exactly \([-1/2,1]\).

Under the common-response restriction write a type as \((x_1,x_2,y_a,y_b,r_c)\).  The mixture with mass \(1/2\) on each of
\[
(c,a,0,0,1),
\qquad
(a,c,0,0,0)
\]
generates \(p^{\mathrm w}\) and has \(\tau=0\).  Replacing both displayed \(y_b\)-coordinates by one leaves the observed law unchanged and gives \(\tau=1\).

For optimality, define
\[
E_1=\{X_1=c,R_c=1\},
\qquad
E_2=\{X_2=c,R_c=0\}.
\]
The observed off-target cells imply \(P(E_1)=P(E_2)=1/2\).  Because \(R_c\) is binary and common across arms, \(E_1\cap E_2=\varnothing\); hence \(E_1\cup E_2\) has probability one.  On \(E_1\), the zero arm-2 probability of exposure \(b\) and the zero cell \(p^{\mathrm w}_2(c,1)\) imply \(X_2=a\) almost surely.  The zero cell \(p^{\mathrm w}_2(a,1)\), together with target consistency, then implies \(Y_a=0\) on \(E_1\).  Symmetrically, on \(E_2\), the zero arm-1 probability of exposure \(b\) and the zero cell \(p^{\mathrm w}_1(c,0)\) imply \(X_1=a\), and \(p^{\mathrm w}_1(a,1)=0\) implies \(Y_a=0\).  Therefore \(Y_a=0\) almost surely.  Binary support now gives
\[
0\le\tau=E(Y_b)\le1.
\]
The two common-response mixtures attain both endpoints while leaving the observed law fixed.  Convex mixtures attain every intermediate value, so the common-response sharp interval is \([0,1]\).  Its strict difference from \([-1/2,1]\) proves that saturated arm-specific off-target responses are a genuine relaxation of the fully excluded common-response benchmark.
